\documentclass[11pt]{article}
\usepackage[margin=1in]{geometry}
\usepackage{amsmath,amssymb,mathtools}
\usepackage{booktabs,longtable,array}
\usepackage{enumitem}
\usepackage{xcolor}
\usepackage{listings}
\usepackage{tikz}
\usepackage{hyperref}
\usepackage{microtype}
\usetikzlibrary{arrows.meta,positioning,fit,shapes.geometric,calc}
\hypersetup{colorlinks=true,linkcolor=black,urlcolor=blue}
\lstset{basicstyle=\ttfamily\small,frame=single,columns=fullflexible,keepspaces=true,showstringspaces=false,breaklines=true}
\newcommand{\R}{\mathbb{R}}
\newcommand{\Qdom}{\mathcal{Q}}
\newcommand{\epsv}{\boldsymbol{\varepsilon}}
\newcommand{\qv}{\mathbf{q}}
\newcommand{\pv}{\mathbf{p}}
\newcommand{\codepath}[1]{\nolinkurl{#1}}
\title{\textbf{AutoIA: From Mathematical Theory to Source Code}\\
Implementation Map Through Version 0.21.0.dev1}
\author{}
\date{August 2026}
\begin{document}
\maketitle

\begin{abstract}
This document connects the mathematical specification of AutoIA to the concrete
Python implementation in version 0.21.0.dev1.  It is intended for a researcher or
developer who understands the enclosure theory and wants to know exactly where
that theory is realized in the source tree.  The organizing principle is not a
line-by-line code listing.  Instead, each mathematical invariant, theorem, or
algorithm is mapped to the production module that proposes it, the module that
certifies it, the representation object that carries it, and the regression or
audit tests that protect it.

A recurring distinction is made between \emph{proposal code} and \emph{proof
code}.  Proposal code is allowed to use ordinary NumPy floating-point operations
because a later proof validates the stored proposal.  Proof code is either in the
trusted binary64 layer or is a composition of trusted interval operations.  This
document highlights that boundary explicitly so future changes do not accidentally
move an unverified calculation into the trusted path.
\end{abstract}

\tableofcontents
\newpage

%===============================================================================
\section{Repository architecture at version 0.21.0.dev1}
%===============================================================================

The production package is organized into layers with deliberately restricted
responsibilities.  The principal directories are:

\begin{longtable}{p{0.29\textwidth}p{0.64\textwidth}}
\toprule
Production location & Responsibility \\
\midrule
\endhead
\codepath{autoia/api.py} & Stable public one-box/subdivision API, input validation, typed exception translation, and user-model guarding. \\
\codepath{autoia/session.py} & Prepared repeated-query sessions, certificate cache, rigorous certificate restriction, serial/thread/process many-box execution, and streaming APIs. \\
\codepath{autoia/trusted/} & Proof-critical outward-rounded binary64 interval arithmetic, runtime environment audit, rigorous reductions, and trusted scalar elementary-function kernels. \\
\codepath{autoia/polynomial/} & Representation-independent sparse polynomial algebra, monomial/range logic, Bernstein transforms, and rigorous coordinate recentering. \\
\codepath{autoia/enclosures/} & Scalar validated representations and representation-level arithmetic: interval, polynomial, shared-generator, piecewise, power, elementary lifting, and runtime contexts. \\
\codepath{autoia/semantic/} & Local NumPy/SciPy dispatch and semantic linear-solve certificates. \\
\codepath{autoia/controller/} & Core execution orchestration, final range policy, subdivision, split scoring, path-aware branching, and diagnostics. \\
\codepath{autoia/elementary/} & Generic Taylor coefficient/remainder machinery above the trusted scalar kernels. \\
\codepath{autoia/special.py} & Public error-function and normal PDF/CDF convenience functions. \\
\codepath{autoia/results.py} & Public result aliases/facade. \\
\codepath{autoia/policies.py} & Public range/subdivision/branch policy classes with typed validation. \\
\codepath{autoia/exceptions.py} & Frozen public exception taxonomy. \\
\bottomrule
\end{longtable}

The low-level dependency direction is intentionally one-way.  In particular,
\codepath{autoia/trusted/} must not import controller, semantic, enclosure, or API
layers.  Boundary tests enforce this architectural rule.

%===============================================================================
\section{One-box call path from public API to proof result}
%===============================================================================

A normal user call is
\begin{lstlisting}[language=Python]
result = autoia(model, parameter_intervals,
                order=4,
                representation="polynomial")
\end{lstlisting}
The production flow is:

\begin{enumerate}[leftmargin=*]
\item \codepath{autoia/api.py}: \texttt{autoia(...)} validates public inputs and
creates/validates the backend and policies.
\item \codepath{autoia/api.py}: \texttt{\_guard\_model(...)} wraps the user model
so user exceptions and AutoIA proof failures are classified separately.
\item \codepath{autoia/controller/execution.py}: the controller constructs
rigorous centered input domains and initial representation objects.
\item \codepath{autoia/semantic/dispatch.py}: the user's function executes in a
private shadow namespace whose NumPy/math/SciPy references can route supported
operations to validated semantics without mutating global modules.
\item Scalar operations are performed by \texttt{ValidatedValue} and its concrete
representation classes under \codepath{autoia/enclosures/}.
\item Recognized \texttt{np.linalg.solve} calls are handled by
\codepath{autoia/semantic/linear_solve.py}.
\item If a Python scalar predicate cannot be decided, path handling under
\codepath{autoia/controller/branching.py} re-executes feasible paths.
\item Final representation ranges are chosen by
\codepath{autoia/controller/ranging.py} and the polynomial range engines under
\codepath{autoia/polynomial/ranging.py}.
\item \codepath{autoia/controller/execution.py} constructs an
\texttt{AutoIAResult}; path-aware execution constructs an
\texttt{AutoIAPathwiseResult}.
\end{enumerate}

The important engineering point is that the user program is still executed as
ordinary Python.  AutoIA changes the objects and selected library references seen
by that program, rather than translating the complete function into a separate
symbolic language.

%===============================================================================
\section{Public input intervals and rigorous centering}
%===============================================================================

The mathematical input contract is
\[
  p_i\in[\underline p_i,\overline p_i].
\]
The public interval class is exposed from \codepath{autoia/intervals.py}, which
facades the trusted interval representation.

The centered-domain algorithm
\[
  q_i=p_i-p_{c,i}\in[-h_i,h_i]
\]
is implemented by
\begin{center}
\codepath{autoia/controller/execution.py}: \texttt{rigorous\_center\_and\_halfwidth}.
\end{center}
The center is an ordinary finite binary64 proposal.  The difference between the
user interval and this center is computed through the trusted backend, and
\(h_i\) is chosen from the resulting rigorous interval.  This is where the known
symmetric-domain boundary enlargement enters the current design.

Initial uncertain parameter objects are constructed in
\texttt{autoia\_single(...)} in the same module.  Depending on the requested
representation, it calls:
\begin{itemize}[leftmargin=*]
\item \texttt{PolynomialEnclosure.parameter(...)};
\item \texttt{SharedGeneratorPolynomialEnclosure.parameter(...)}; or
\item \texttt{IntervalEnclosure(...)}.
\end{itemize}

%===============================================================================
\section{Representation-independent value dispatch}
%===============================================================================

The abstract mathematical invariant
\[
  z(\qv)\in\Gamma_E(\qv)
\]
is represented in code by a \texttt{ValidatedValue} wrapper and a concrete
representation object.

The representation-independent wrapper is in
\begin{center}
\codepath{autoia/enclosures/value.py}.
\end{center}
It provides arithmetic operator hooks and delegates to a promoted concrete
representation.  The promotion priority is strategy, not proof.  Soundness comes
from each target representation's \texttt{convert\_from(...)} implementation.

The currently important concrete scalar families are:

\begin{longtable}{p{0.30\textwidth}p{0.30\textwidth}p{0.32\textwidth}}
\toprule
Mathematical family & Python class & Production file \\
\midrule
\endhead
\(P(\qv)+R\) & \texttt{ParameterizedEnclosure} / \texttt{PolynomialEnclosure} & \codepath{autoia/enclosures/polynomial.py} \\
Single interval \(I\) & \texttt{IntervalEnclosure} & \codepath{autoia/enclosures/interval.py} \\
\(P(\qv)+G\epsv+R\) & \texttt{SharedGeneratorPolynomialEnclosure} & \codepath{autoia/enclosures/generator.py} \\
Finite path union & \texttt{PiecewiseEnclosure} & \codepath{autoia/enclosures/piecewise.py} \\
\bottomrule
\end{longtable}

The generator representation advertises a higher promotion priority than the
plain polynomial representation so shared provenance is retained when sound
conversion is available.

%===============================================================================
\section{Polynomial algebra: equations to functions}
%===============================================================================

Sparse point-coefficient polynomial algebra lives in
\codepath{autoia/polynomial/algebra.py}.

\subsection{Multi-indices}

The mathematical objects
\[
  \alpha=(\alpha_1,\ldots,\alpha_m),
  \qquad
  |\alpha|=\sum_i\alpha_i
\]
are Python tuples of integers.  Important helper functions include:
\begin{itemize}[leftmargin=*]
\item \texttt{zero\_alpha(m)};
\item \texttt{basis\_alpha(m,i)};
\item \texttt{total\_degree(alpha)};
\item \texttt{multi\_indices\_exact\_degree(m,degree)}.
\end{itemize}
The exact-degree multi-index collections are cached because they are purely
combinatorial and independent of numerical values.

\subsection{Monomial magnitude}

The proof quantity
\[
  H_\alpha\ge\max_{\qv\in\Qdom}|\qv^\alpha|
\]
is computed by \texttt{monomial\_magnitude\_upper}.  Its parity-aware interval
version is \texttt{monomial\_range\_symmetric}.  Backend-local caching avoids
recomputing the same proof bound for repeated \((\alpha,h)\) combinations.

\subsection{Addition}

The rigorous point-coefficient addition rule is implemented by
\texttt{rigorous\_poly\_add\_point}.  The function returns both the stored point
polynomial proposal and an interval enclosing coefficient-rounding error after
that error has been propagated through monomial ranges.

\subsection{Truncated multiplication}

The mathematical coefficient convolution
\[
  C_\gamma=\sum_{\alpha+\beta=\gamma}c_\alpha d_\beta
\]
with a retained degree cutoff is implemented by
\texttt{rigorous\_poly\_mul\_truncated}.  Sparse convolution structure is cached
by \texttt{\_polynomial\_product\_plan}.  The numerical coefficients themselves
are not cached.  Every kept coefficient is stored as a point proposal with
rigorous coefficient error, and discarded high-degree contributions are ranged
into an interval remainder.

\subsection{Reciprocal proposal}

\texttt{build\_reciprocal\_predictor\_point} constructs a point polynomial
proposal for a reciprocal.  The representation-level reciprocal certificate is
implemented above this helper in \codepath{autoia/enclosures/polynomial.py}; it
uses the residual identity rather than trusting the proposal.

%===============================================================================
\section{Polynomial range methods in code}
%===============================================================================

All principal polynomial range engines are in
\codepath{autoia/polynomial/ranging.py}.

\begin{longtable}{p{0.37\textwidth}p{0.55\textwidth}}
\toprule
Function & Mathematical role \\
\midrule
\endhead
\texttt{range\_termwise\_scalar} & Sum rigorous monomial intervals. \\
\texttt{derivative\_range\_termwise} & Rigorous range of \(\partial P/\partial q_i\). \\
\texttt{range\_monotonicity\_scalar} & Detect sign-definite derivatives and fix monotone coordinates at extrema. \\
\texttt{power\_coeff\_unit\_box\_interval} & Convert power coefficients toward the full Bernstein representation on a unit box. \\
\texttt{range\_bernstein\_scalar} & Rigorous Bernstein coefficient hull. \\
\texttt{power\_coeff\_reduced\_unit\_box\_interval} & Bernstein conversion after monotone coordinates have been fixed. \\
\texttt{range\_monotonicity\_reduced\_bernstein\_scalar} & Combined structural range method used by the adaptive policy. \\
\texttt{range\_polynomial} & Public/internal range-method selector and portfolio. \\
\bottomrule
\end{longtable}

Final output policy is selected by \codepath{autoia/controller/ranging.py}.  The
controller layer decides which sound range result to request; the polynomial
module performs the actual mathematical enclosure.

%===============================================================================
\section{Trusted binary64 layer: the proof boundary}
%===============================================================================

The trusted layer is the most important place to maintain architectural discipline.
Its production files are:

\begin{itemize}[leftmargin=*]
\item \codepath{autoia/trusted/environment.py}: runtime floating-point audit;
\item \codepath{autoia/trusted/binary64.py}: scalar outward-rounded arithmetic;
\item \codepath{autoia/trusted/interval.py}: trusted interval type and elementary interval operations;
\item \codepath{autoia/trusted/reductions.py}: rigorous vector/matrix reductions;
\item \codepath{autoia/trusted/elementary.py}: trusted scalar elementary and special functions.
\end{itemize}

The public backend facade is exposed as \texttt{RigorousBinary64Backend} from
\codepath{autoia/backends.py}.  The underlying implementation is imported from
the trusted package.

\subsection{Long dot products}

The rigorous dot-product contract
\[
  \sum_i a_i b_i\in I
\]
is implemented by \texttt{rigorous\_dot\_point}.  It selects either the
historical scalar reference implementation or the vectorized outward pairwise
implementation according to the backend's \texttt{reduction\_method} and term
threshold.

Step 9b extends the same proof pattern to
\begin{itemize}[leftmargin=*]
\item \texttt{rigorous\_dot\_point\_interval};
\item \texttt{rigorous\_dot\_interval};
\item point/interval matrix-vector products;
\item point/interval matrix-matrix products;
\item nonnegative upper-bound reductions.
\end{itemize}
The batched matrix kernels are intentionally written as collections of explicit
outward product intervals and pairwise interval sums.  They are not wrappers
around an unverified BLAS reduction.

%===============================================================================
\section{Elementary functions: trusted scalar kernel plus generic Taylor lifting}
%===============================================================================

The implementation is split deliberately.

\paragraph{Trusted scalar layer.}
\codepath{autoia/trusted/elementary.py} implements rigorous scalar/interval
kernels for \(\sqrt{\cdot}\), \(\exp\), \(\log\), \(\sin\), \(\cos\),
\(\operatorname{erf}\), normal PDF/CDF, and constant powers.  Important helpers
include \texttt{pi\_interval}, \texttt{ln2\_interval}, and the point/interval
variants of every supported function.

\paragraph{Generic polynomial lifting.}
\codepath{autoia/elementary/taylor.py} builds Taylor coefficient intervals and
remainder radii.  \texttt{taylor\_coefficients} and
\texttt{taylor\_remainder\_radius} contain the mathematical function-specific
recurrences and tail bounds; \texttt{build\_taylor\_data} packages them for the
enclosure layer.

\paragraph{Representation dispatch.}
\codepath{autoia/enclosures/elementary.py} applies the Taylor data to polynomial,
interval, and generator-valued arguments.  This division keeps transcendental
proof facts out of user dispatch code.

%===============================================================================
\section{Power semantics in code}
%===============================================================================

Constant exponent handling is centralized in
\codepath{autoia/enclosures/power.py} and trusted support in
\codepath{autoia/trusted/elementary.py}.  Validated arithmetic routes
\texttt{x ** p}, \texttt{pow(x,p)}, and supported NumPy/math power calls to this
layer.  The trusted functions \texttt{real\_power\_point\_interval} and
\texttt{real\_power\_interval} enforce the real-domain restrictions described in
the mathematical specification.

%===============================================================================
\section{Local NumPy/math/SciPy semantic dispatch}
%===============================================================================

Global monkeypatching is deliberately avoided.  The key code is
\codepath{autoia/semantic/dispatch.py}.

\subsection{Prepared shadow function}

\texttt{prepare\_model\_for\_dispatch} and
\texttt{prepare\_dispatch\_model} clone the user function with a private shadow
global namespace.  Same-module helper functions can also be cloned into that
namespace.  NumPy and \texttt{math} names are replaced by proxies that route only
supported validated calls; the real modules remain unchanged.

The class \texttt{PreparedDispatchModel} is the static object reused by
\texttt{PreparedAutoIA} in Step 9c.

\subsection{NumPy proxy}

\texttt{\_NumpyProxy} and \texttt{\_LinalgProxy} provide local dispatch for
array constructors, elementary operations, norms, and \texttt{linalg.solve}.
Object-array construction is used so ordinary NumPy indexing/broadcasting invokes
validated scalar arithmetic.

\subsection{SciPy compatibility}

\codepath{autoia/semantic/scipy_dispatch.py} recognizes supported public SciPy
surfaces such as \texttt{scipy.special.erf}, \texttt{ndtr}, legacy
\texttt{scipy.stats.norm}, frozen normal distributions, and the newer
\texttt{scipy.stats.Normal} API when present.  SciPy syntax is a front-end only;
the proof routes to AutoIA's own rigorous normal/error-function kernels.

%===============================================================================
\section{Semantic linear solve: theory-to-code map}
%===============================================================================

The main implementation is
\codepath{autoia/semantic/linear_solve.py}.  It contains both proposal builders
and certificate routines.  This file is intentionally larger than most modules
because it holds several no-regression solve paths.

\subsection{General polynomial predictor}

The recurrence
\[
  A_0x_\alpha=b_\alpha-
  \sum_{\substack{\beta\ne0\\\beta\le\alpha}}
  A_\beta x_{\alpha-\beta}
\]
is implemented by
\texttt{build\_solution\_polynomials\_general}.  Coefficients of the same total
degree are solved together as a multiple-right-hand-side NumPy solve so one
reference factorization can serve a whole degree.  This is proposal-only
performance work from Step 9.

\subsection{Residual coefficient construction}

\texttt{residual\_coefficient\_intervals\_general} constructs rigorous interval
coefficients for
\[
  r=b-AX.
\]
\texttt{range\_interval\_coefficient\_poly} and related functions rigorously
range the resulting interval-coefficient polynomials.

\subsection{Preconditioner portfolio}

\texttt{polynomial\_preconditioner\_portfolio} generates center and matrix-family
representatives.  The output candidates are ordinary point matrices.  Candidate
quality is not trusted.

\subsection{Defect and contraction proof}

\texttt{preconditioned\_defect\_bounds\_general} establishes a rigorous matrix
bound for
\[
  |I-CA|.
\]
\texttt{evaluate\_general\_polynomial\_candidate} combines residual and defect
bounds, checks \(\rho<1\), constructs the scalar/componentwise correction, and
returns the certified candidate.

The older affine path remains separately implemented for historical bit-for-bit
regression.  \texttt{polynomial\_semantic\_solve\_general} selects the general
path only when required.

\subsection{Generator-aware solve}

The Step-7b path begins around
\texttt{generator\_semantic\_solve}.  Important functions include:
\begin{itemize}[leftmargin=*]
\item \texttt{\_generator\_sensitivity\_proposal};
\item \texttt{\_generator\_predictor\_vector};
\item \texttt{\_generator\_residual\_vector};
\item \texttt{\_generator\_defect\_bounds};
\item \texttt{\_evaluate\_generator\_solve\_candidate}.
\end{itemize}
These functions retain incoming generator identifiers in the explicit solution
predictor and prove the remaining error by the same contraction theorem over the
joint physical/generator domain.

\subsection{Dynamic solve context}

\texttt{SolveContext} and \texttt{GenericSolveContext} provide the run-local
semantic solve dispatcher.  The local dispatch mechanism passes the current
context explicitly/context-locally rather than replacing global
\texttt{np.linalg.solve}.

%===============================================================================
\section{Path-aware control flow in code}
%===============================================================================

Path logic lives principally in \codepath{autoia/controller/branching.py} and
supporting predicate objects in \codepath{autoia/enclosures/branching.py}.

The internal \texttt{\_PathState} tracks the current physical box, forced
Boolean outcomes already chosen during path enumeration, split depth, and trace
information.  \texttt{autoia\_pathwise} is the path-aware execution controller.

For exact affine one-variable predicates, \texttt{\_exact\_single\_variable\_root}
tries to recover a direct physical cut location.  Otherwise
\texttt{\_predicate\_split\_direction} proposes a geometric split direction.
\texttt{\_make\_piecewise\_result} assembles completed feasible paths into the
public piecewise result.

Internal branch-control signals derive from \texttt{BaseException} rather than
ordinary \texttt{Exception}; this prevents a user's broad
\texttt{except Exception:} from swallowing AutoIA's internal request to re-run a
path under a forced predicate outcome.

%===============================================================================
\section{Adaptive subdivision and cheap split scoring}
%===============================================================================

Prescribed and adaptive subdivision are implemented by
\codepath{autoia/controller/subdivision.py}.  The practical cheap score is
separated into \codepath{autoia/controller/split_scoring.py}.

The score's polynomial re-centering is intentionally a cheap floating prediction,
not a proof.  Functions such as \texttt{translate\_polynomial\_one\_direction},
\texttt{derivative\_range\_termwise\_float}, and
\texttt{score\_split\_direction} estimate the structural effect of a hypothetical
half split.  \texttt{cheap\_score\_key} implements the lexicographic ordering.

After a direction is selected, \codepath{autoia/controller/subdivision.py}
actually certifies the chosen two children and measures realized gain before
accepting the split.  The historical reference mode remains available through
\texttt{split\_method="lookahead"}.

%===============================================================================
\section{Prepared sessions and certificate restriction in code}
%===============================================================================

Step 9c is implemented primarily in \codepath{autoia/session.py} and
\codepath{autoia/polynomial/recenter.py}.

\subsection{Public preparation object}

\texttt{prepare(model,...)} returns \texttt{PreparedAutoIA}.  Its constructor:
\begin{enumerate}[leftmargin=*]
\item validates the fixed model/configuration;
\item creates or accepts one \texttt{RigorousBinary64Backend};
\item creates one \texttt{CertificateCache};
\item wraps the user model through \texttt{\_guard\_model}; and
\item calls \texttt{prepare\_dispatch\_model} exactly once.
\end{enumerate}
The session then exposes \texttt{certify}, \texttt{certify\_many},
\texttt{certify\_many\_bounds}, \texttt{iter\_certify\_many}, and
\texttt{iter\_bounds\_many}.

\subsection{Exact cache}

\texttt{CertificateCache} is a bounded LRU cache keyed by exact stored endpoint
tuples.  \texttt{get\_exact} is an O(1) dictionary lookup.  An exact cache hit
returns a shallow result wrapper so per-request \texttt{cache\_info} metadata
does not mutate a result object previously returned to the user.

\subsection{Containment cache}

\texttt{get\_smallest\_container} scans the bounded cache for parent boxes that
contain the query and chooses the smallest by log-volume.  This choice is a
heuristic among already sound candidate parents; no proof depends on it.

\subsection{Restriction function}

\texttt{\_restrict\_cached\_result} performs the rigorous subbox restriction.
For every new input interval it calls
\texttt{rigorous\_center\_and\_halfwidth}.  For polynomial/generator
representations it constructs the actual outward symmetric physical envelope and
requires that envelope to lie inside the cached parent box.

The exact real center shift is enclosed by trusted subtraction:
\begin{lstlisting}[language=Python]
backend.sub(point_interval(c_new), point_interval(c_old))
\end{lstlisting}
The resulting intervals are passed to
\texttt{rigorous\_recenter\_point\_polynomial} in
\codepath{autoia/polynomial/recenter.py}.

That function implements the multivariate binomial coordinate transform.  For
every exact transformed coefficient it simultaneously computes:
\begin{itemize}[leftmargin=*]
\item a rigorous interval using trusted operations and rigorous shift intervals;
\item an ordinary binary64 point proposal using the stored shift proposal.
\end{itemize}
The difference is multiplied by a rigorous monomial range over the new half-width
vector and accumulated in \texttt{extra\_remainder}.  The new representation is
then final-ranged using the normal adaptive range policy.

\subsection{Interval representation restriction}

An interval-only cached result needs no polynomial coordinate transform.  The
same interval enclosure remains valid on any physical subbox.  This can eliminate
model execution but generally does not sharpen the output interval unless a fresh
proof is performed.

%===============================================================================
\section{Many-box CPU execution in code}
%===============================================================================

\texttt{PreparedAutoIA.certify\_many} normalizes a finite collection of boxes and
supports three execution modes.

\paragraph{Serial.}
The session calls \texttt{certify} in input order and shares one cache/backend.

\paragraph{Thread.}
\texttt{ThreadPoolExecutor} runs chunks against the same prepared session.  Locks
protect session statistics/cache bookkeeping, while proof runtime state is
context-local.  The current Python-object workload is not guaranteed to benefit
from threads.

\paragraph{Process.}
\texttt{ProcessPoolExecutor} with the \texttt{spawn} context uses
\texttt{\_process\_initializer}.  Each worker constructs one independent trusted
backend and one \texttt{PreparedAutoIA} session, then reuses it for all assigned
chunks.  The model must be pickleable, normally meaning it is a module-level
function.

\texttt{iter\_certify\_many} uses a bounded deque of pending futures so only a
small number of chunks are in flight.  This is the API intended for very large
streams.  \texttt{iter\_bounds\_many} discards full result objects after yielding
copies of lower/upper arrays.

%===============================================================================
\section{Result objects and mathematical diagnostics}
%===============================================================================

The primary result dataclass is defined in
\codepath{autoia/controller/execution.py} as \texttt{AutoIAResult}.  It stores:
\begin{itemize}[leftmargin=*]
\item flat \texttt{lower} and \texttt{upper} arrays;
\item structured validated \texttt{objects};
\item backend name and \texttt{machine\_rigorous} flag;
\item semantic certificates;
\item requested representation;
\item range policy and final-range diagnostics;
\item output shape.
\end{itemize}
Convenience properties reconstruct shaped outputs and expose generator IDs/matrix
when present.

Path-aware results are implemented in
\codepath{autoia/controller/branching.py} and exposed through
\codepath{autoia/results.py}.  They retain per-branch boxes, decisions, traces,
and piecewise output enclosures in addition to global lower/upper hulls.

Prepared-session results receive a lightweight \texttt{cache\_info} attribute
identifying whether the query was a fresh certificate, an exact cache hit, or a
restricted parent certificate.

%===============================================================================
\section{Typed failure semantics and where failures are translated}
%===============================================================================

Public exception classes are defined in \codepath{autoia/exceptions.py}.  The
important categories are:
\begin{itemize}[leftmargin=*]
\item \texttt{InvalidInputError};
\item \texttt{UnsupportedOperationError};
\item \texttt{CertificationError};
\item \texttt{ResourceLimitError};
\item \texttt{NumericalLimitError};
\item \texttt{ModelExecutionError};
\item \texttt{InternalAutoIAError}.
\end{itemize}

\codepath{autoia/api.py}: \texttt{\_translate\_internal\_exception} maps lower
level failures to this taxonomy.  The distinction between a failed proof and an
unsupported operation is important: a contraction failure means the requested
claim was not proved on that box, not that the model syntax is unsupported.

%===============================================================================
\section{Testing and audit map}
%===============================================================================

The production source is protected by several kinds of tests.

\subsection{Historical bit-for-bit regression}

The retired Version-0.11 reference implementation lives under
\codepath{tests/support/oracle_v11.py}.  Tests such as
\codepath{tests/test_reference_regression.py},
\codepath{tests/test_semantic_solve_regression.py}, and trusted-kernel regression
tests protect established arithmetic paths from accidental changes.

\subsection{Exact-rational trusted-kernel tests}

\codepath{tests/test_scalar_binary64_exact.py},
\codepath{tests/test_reductions_exact.py}, and Step-9/9b performance tests compare
proof-critical binary64 enclosures to exact rational calculations.

\subsection{High-precision elementary tests}

Trusted transcendental functions are checked against independent high-precision
references in \codepath{tests/test_trusted_elementary.py}, power tests, and the
Step-3 audit tools.

\subsection{Representation and semantic tests}

Polynomial, generator, solve, branching, SciPy compatibility, array semantics,
and controller tests exercise each representation and semantic boundary.

\subsection{Prepared-session tests}

\codepath{tests/test_prepared_manybox.py} verifies:
\begin{itemize}[leftmargin=*]
\item one-shot versus prepared fresh equivalence;
\item exact-cache model-execution avoidance;
\item per-request cache metadata isolation;
\item interior certificate restriction;
\item fail-closed boundary reuse;
\item generator-ID preservation through restriction;
\item thread order preservation;
\item compact bounds and streaming APIs;
\item process-mode validation for non-pickleable functions.
\end{itemize}
The separate tool \codepath{tools/audit_manybox.py} measures repeated-query
performance, checks sampled truth containment, and exercises process-mode outer
parallelism.

%===============================================================================
\section{How to modify AutoIA without breaking the proof architecture}
%===============================================================================

This section gives implementation recipes tied directly to the mathematical
contracts.

\subsection{Adding a new scalar elementary function}

A mathematically smooth scalar function should normally be added in four places:
\begin{enumerate}[leftmargin=*]
\item add a trusted scalar/interval kernel or a trusted domain/range primitive in
\codepath{autoia/trusted/elementary.py};
\item add Taylor coefficient/remainder formulas in
\codepath{autoia/elementary/taylor.py};
\item connect representation-level lifting in
\codepath{autoia/enclosures/elementary.py};
\item add local NumPy/math/SciPy dispatch only for actual public syntax that
users are expected to write.
\end{enumerate}
The trusted layer must not call unverified platform transcendental functions as
proof facts unless a separate rounding guarantee is established.

\subsection{Adding a new semantic numerical primitive}

For an operation such as a future nonlinear root solver or eigensystem:
\begin{enumerate}[leftmargin=*]
\item define the exact mathematical semantics being certified;
\item permit arbitrary ordinary-floating proposal calculations;
\item derive an independent residual/fixed-point/inclusion certificate;
\item place the semantic proof under \codepath{autoia/semantic/};
\item route recognized syntax through local dispatch without global monkeypatching;
\item add generic fallback behavior or a typed unsupported failure when the proof
cannot apply.
\end{enumerate}
The semantic primitive should not be accepted merely because an approximate
library routine returned a plausible result.

\subsection{Adding a new representation}

A new representation must define a concretization and sound arithmetic/conversion
rules.  It should implement the representation contract expected by
\texttt{ValidatedValue}, including a stable \texttt{representation\_name}, range
operation, arithmetic where supported, and inclusion-preserving
\texttt{convert\_from}.  Promotion priority is chosen only after the conversions
are proved sound.

\subsection{Adding a new range method}

Keep the proof implementation under \codepath{autoia/polynomial/} or an analogous
low-level range module.  Controller policy may decide when to use it, but the
method itself should return an interval independently known to contain the exact
represented range.  A portfolio may then intersect it with other rigorous
methods.

\subsection{Changing the prepared-session cache}

The public \texttt{PreparedAutoIA} API should not require a specific cache data
structure.  The current bounded LRU scan can later be replaced by a spatial index.
A new index may alter parent-selection cost and reuse frequency, but must never
skip the explicit proof that the query's actual represented centered domain lies
inside the parent's certified domain before polynomial restriction.

\subsection{Adding a GPU backend}

The current Python object representation should not simply be copied to a GPU.
The preferred architecture is a numeric structure-of-arrays backend behind the
already frozen prepare/certify-many public API.  GPU proposal kernels can be
introduced first because their outputs can be revalidated by the existing CPU
trusted proof layer.  A future proof-producing GPU backend would require its own
audited directed-rounding primitives and independent exact-reference tests.

%===============================================================================
\section{Proposal-versus-proof checklist for code review}
%===============================================================================

Before accepting a modification, ask:
\begin{enumerate}[leftmargin=*]
\item What exact mathematical set or equation is being certified?
\item Which computed quantities are only proposals?
\item Where is the exact residual, interval extension, or domain condition checked?
\item Does any proof-critical calculation use ordinary round-to-nearest arithmetic
without an enclosing error calculation?
\item If a representation loses structure, is the conversion explicitly
inclusion-preserving?
\item If correlation is shared, what proves that two terms refer to the same
auxiliary variable?
\item If a path is discarded, what proves it is infeasible?
\item If a cached parent certificate is reused, what proves the new represented
domain is contained in the parent domain?
\item If performance batching changes arithmetic order, is every batched reduction
still a composition of rigorous interval operations?
\item Is there an exact/high-precision/regression test that would fail if the new
proof boundary were accidentally weakened?
\end{enumerate}

%===============================================================================
\section{Compact source-map summary}
%===============================================================================

\begin{longtable}{p{0.31\textwidth}p{0.61\textwidth}}
\toprule
Mathematical idea & Main implementation location \\
\midrule
\endhead
Input centering \(p\mapsto q\) & \codepath{autoia/controller/execution.py} \\
Outward scalar interval arithmetic & \codepath{autoia/trusted/binary64.py}, \codepath{autoia/trusted/interval.py} \\
Rigorous dot/matrix reductions & \codepath{autoia/trusted/reductions.py} \\
Polynomial algebra and truncation & \codepath{autoia/polynomial/algebra.py} \\
Termwise/monotonicity/Bernstein range & \codepath{autoia/polynomial/ranging.py} \\
Polynomial coordinate restriction/recentering & \codepath{autoia/polynomial/recenter.py} \\
\(P+R\) representation & \codepath{autoia/enclosures/polynomial.py} \\
\(P+G\varepsilon+R\) representation & \codepath{autoia/enclosures/generator.py} \\
Piecewise representation & \codepath{autoia/enclosures/piecewise.py} \\
Taylor lifting & \codepath{autoia/elementary/taylor.py}, \codepath{autoia/enclosures/elementary.py} \\
Trusted transcendentals & \codepath{autoia/trusted/elementary.py} \\
Semantic linear solve & \codepath{autoia/semantic/linear_solve.py} \\
NumPy/math dispatch & \codepath{autoia/semantic/dispatch.py} \\
SciPy normal/error-function syntax & \codepath{autoia/semantic/scipy_dispatch.py} \\
Path-aware execution & \codepath{autoia/controller/branching.py} \\
Adaptive split score & \codepath{autoia/controller/split_scoring.py} \\
Subdivision orchestration & \codepath{autoia/controller/subdivision.py} \\
Prepared/many-box execution & \codepath{autoia/session.py} \\
Public error taxonomy & \codepath{autoia/exceptions.py}, translated in \codepath{autoia/api.py} \\
\bottomrule
\end{longtable}

The main design principle behind this map is that the mathematical proof should
remain visible in the source architecture.  Proposal code may evolve rapidly for
performance, but every accepted certificate should still be traceable to a small
number of explicitly audited enclosure operations and documented mathematical
identities.

\end{document}
